% Generated by SIGMA; compile with Tectonic or another TeX engine.
\documentclass[11pt]{article}
\usepackage[T1]{fontenc}
\usepackage[utf8]{inputenc}
\usepackage[margin=1in]{geometry}
\usepackage{amsmath,amssymb,booktabs,longtable}
\usepackage{array}
\usepackage{enumitem}
\setlist{nosep,leftmargin=*}
\usepackage{microtype}
\usepackage{titlesec}
\usepackage{xcolor}
\definecolor{blue}{HTML}{24485E}
\titleformat{\section}{\Large\bfseries\sffamily\color{blue}}{\thesection}{0.8em}{}
\titleformat{\subsection}{\large\bfseries\sffamily}{\thesubsection}{0.7em}{}
\titleformat{\subsubsection}{\normalsize\bfseries\sffamily\color{blue}}{\thesubsubsection}{0.6em}{}
\setlength{\parskip}{0.55em}
\setlength{\parindent}{0pt}
\title{Maximum Regions Formed by Lines v1}
\author{SIGMA Research Record}
\date{\today}
\begin{document}
\maketitle
\begin{quote}\small In the future we plan on using AI to produce a typical academic paper based on the results of the research.\end{quote}
\section{Project}
\textbf{Interpretation:} Agent-generated findings are provisional unless explicitly verified; reviewer agreement is not proof.\par
\textbf{Research problem:} For each integer n \textgreater{}= 0, determine the maximum number of regions into which n distinct straight lines can divide the Euclidean plane. Lines may be placed arbitrarily. Derive a formula for the maximum, prove that no arrangement can exceed it, and show that an arrangement attaining it exists for every n. Include the n = 0 case and check small values. Do not assume a particular proof method.\par
\textbf{Project ID:} 5731fd49-e4fd-47c3-9d3e-1691e30d6ea1\par
\textbf{Research type:} Mathematics\par
\textbf{Difficulty:} Easy\par
\section{Strategies}
\subsection{Incremental region counting with an explicit maximizing arrangement (saved)}
Define a region as a connected component of the plane after removing the lines. With no lines, the whole plane is one region, so R(0) = 1.

For the upper bound, insert the lines one at a time. The kth line meets the preceding k − 1 lines in at most k − 1 distinct points. Those points divide the new line into at most k open intervals or rays. Each such piece lies in an existing region. Every existing region is convex: it is the intersection of the open half-planes specifying which side of each existing line its points occupy. Consequently, a region cannot contain two separated pieces of the new line. A region met by a piece is split into exactly two regions, since it is open and convex and has points on both sides of the new line. Thus the kth line adds at most k regions. Summing gives R(n) ≤ 1 + ∑(k = 1 to n) k = 1 + n(n + 1)/2.

For attainment, take the n lines L\_k: y = kx + k², for k = 1, …, n. When L\_k is inserted, its intersection with L\_i, where i \textless{} k, has x-coordinate −(i + k). These k − 1 coordinates are distinct. Hence L\_k is divided into exactly k pieces by the earlier lines. By the convex-region argument above, these pieces lie in distinct existing regions, each of which is split in two. Therefore L\_k adds exactly k regions. This arrangement attains 1 + n(n + 1)/2 for every positive n; the empty arrangement attains it for n = 0. Together with the upper bound, this proves the formula.\par
\subsection{Separate the bound from attainment (saved)}
Prove a new line crosses at most n-1 distinct old intersection points, then show its pieces split distinct convex regions. Independently verify an explicit family with no parallel pairs or triple concurrence. Include n=0 and n=1. Evaluator-authored plan, not a saved session result.\par
\section{Research sessions}
No research sessions are recorded.\par
\section{Agentic runs}
No agentic runs are recorded.\par
\end{document}
